% ======================================================================
\chapter[Фотоны. Фотоэффект. Эффект Комптона]{Фотоны. Фотоэффект. Эффект Комптона}
\label{t:01}
% ======================================================================

\begin{qbox}
\textbf{Билет 1.} Квантовая природа электромагнитного поля (фотоны). Уравнение Эйнштейна
и объяснение фотоэффекта. Эффект Комптона, основные соотношения. Комптоновская длина волны.
\end{qbox}

\section{Содержание билета}

\textbf{Подтемы:}
\begin{enumerate}
  \item фотон как квант электромагнитного поля: энергия, импульс, масса;
  \item экспериментальные закономерности фотоэффекта и их объяснение; уравнение Эйнштейна;
  \item эффект Комптона: опыт, формула для сдвига длины волны, энергия рассеянного фотона и электрона отдачи;
  \item комптоновская длина волны и её физический смысл.
\end{enumerate}

\textbf{Необходимые понятия:} законы сохранения энергии и импульса; релятивистская связь
$E^2=p^2c^2+m^2c^4$; плоская электромагнитная волна, её частота $\omega$ и волновой вектор
$\vect k$ ($k=\omega/c$); интенсивность излучения; работа выхода электрона из металла.

\section{Определения и физические принципы}

\begin{definition}[фотон]
\emph{Фотон} — квант электромагнитного поля: элементарное возбуждение поля, отвечающее
моде (плоской монохроматической волне) с частотой $\omega$ и волновым вектором $\vect k$.
Энергия и импульс фотона
\begin{equation}\label{eq:1:photon}
  E=\hbar\omega=h\nu,\qquad \vect p=\hbar\vect k,\qquad |\vect p|=\frac{\hbar\omega}{c}=\frac{h}{\lambda}.
\end{equation}
Масса покоя фотона равна нулю, он всегда движется со скоростью $c$; спин фотона равен 1,
но из-за поперечности поля возможны лишь два состояния поляризации.
\end{definition}

\emph{Физический смысл.} Энергия электромагнитного поля частоты $\omega$ может меняться только
порциями $\hbar\omega$: поле обменивается энергией и импульсом с веществом дискретно, «поштучно».
Волновые свойства света (интерференция, дифракция) при этом сохраняются: распределение
интенсивности описывает распределение вероятности регистрации фотонов.

\emph{Область применимости.} Корпускулярная картина обязательна, когда существенны отдельные
акты взаимодействия света с атомами и электронами (фотоэффект, эффект Комптона, спектры
излучения, счёт фотонов). Классическая электродинамика справедлива при большом числе фотонов в
каждой моде, когда дискретность $\hbar\omega$ пренебрежимо мала по сравнению с энергией поля.

\begin{definition}[фотоэффект]
\emph{Внешний фотоэффект} — испускание электронов веществом под действием электромагнитного
излучения. \emph{Работа выхода} $A$ — минимальная энергия, которую нужно сообщить электрону,
чтобы удалить его из металла в вакуум. \emph{Красная граница} фотоэффекта — минимальная частота
$\omega_0=A/\hbar$ (максимальная длина волны $\lambda_0=2\pi\hbar c/A$), при которой фотоэффект
ещё возможен. \emph{Задерживающее напряжение} $U_{\text з}$ — наименьшая тормозящая разность
потенциалов, при которой фототок обращается в нуль.
\end{definition}

\begin{definition}[эффект Комптона]
\emph{Эффект Комптона} — увеличение длины волны рентгеновского или гамма-излучения при
рассеянии на свободных (или слабо связанных) электронах. Сдвиг длины волны зависит только
от угла рассеяния $\theta$. \emph{Комптоновская длина волны} частицы массы $m$:
\begin{equation}
  \lambda_{\mathrm C}=\frac{h}{mc},\qquad
  \lambda_{\mathrm C}^{(e)}=\frac{h}{\mel c}=2{,}426\cdot10^{-10}\ \text{см}=2{,}43\pmet .
\end{equation}
Часто используют и приведённую длину $\bar\lambda_{\mathrm C}=\hbar/mc=\lC/2\pi=3{,}86\cdot10^{-11}$~см.
\end{definition}

\emph{Связь понятий.} Фотоэффект доказывает, что энергия поглощается порциями $\hbar\omega$;
эффект Комптона — что фотон несёт также импульс $\hbar\vect k$ и сталкивается с электроном
как частица. Вместе они устанавливают полный набор корпускулярных характеристик~\eqref{eq:1:photon}.

\section{Основная теория}

\subsection{Проблема, приведшая к понятию фотона}

Классическая электродинамика рассматривает свет как непрерывную волну, энергия которой
пропорциональна квадрату амплитуды и может принимать любые значения. К концу XIX века эта
картина столкнулась с двумя трудностями. (1) Теория теплового излучения приводила к
бесконечной полной энергии равновесного излучения («ультрафиолетовая катастрофа»).
М.~Планк (1900) устранил расходимость, предположив, что осциллятор с частотой $\omega$
обменивается энергией с полем порциями $\hbar\omega$. (2) Закономерности фотоэффекта
(Г.~Герц, А.\,Г.~Столетов, Ф.~Ленард) противоречили волновой картине. А.~Эйнштейн (1905)
сделал следующий шаг: квантовано само излучение, свет состоит из квантов с энергией $\hbar\omega$.

\subsection{Энергия, импульс и масса фотона}

Для любой частицы в теории относительности
\begin{equation}\label{eq:1:Ep}
  E^2=p^2c^2+m^2c^4,\qquad \vect v=\frac{\vect pc^2}{E}.
\end{equation}
Частица, движущаяся со скоростью $c$, должна иметь $|\vect p|c=E$, т.\,е. $m=0$. Подставляя
$E=\hbar\omega$, получаем $p=\hbar\omega/c=\hbar k$, а так как направление импульса совпадает с
направлением распространения волны, $\vect p=\hbar\vect k$. Если плотность потока фотонов
равна $j$ (фотонов на единицу площади в единицу времени), то интенсивность $I=j\hbar\omega$, а
давление света при полном поглощении $P=j\hbar\omega/c=I/c$ — тот же результат, что и в
классической теории.

\subsection{Экспериментальные законы фотоэффекта}

Опыт ставится с вакуумным фотоэлементом: освещаемый катод и анод, между которыми задаётся
напряжение $U$. Зависимость фототока $I$ от $U$ (вольт-амперная характеристика) показана на
рис.~\ref{fig:1:photo}. Установлены следующие законы.

\begin{enumerate}
  \item \textbf{Ток насыщения} $I_{\text{нас}}$ (все вылетевшие электроны достигают анода)
        пропорционален интенсивности падающего света при неизменной частоте.
  \item \textbf{Максимальная кинетическая энергия} фотоэлектронов $T_{\max}=eU_{\text з}$
        не зависит от интенсивности света и линейно растёт с частотой.
  \item \textbf{Красная граница:} при $\omega<\omega_0$ фотоэффект отсутствует при сколь угодно
        большой интенсивности; $\omega_0$ определяется материалом катода.
  \item \textbf{Безынерционность:} фототок возникает практически мгновенно
        (запаздывание $<10^{-9}$~с) даже при очень слабом освещении.
\end{enumerate}

\begin{figure}[htbp]
\centering
\begin{tikzpicture}
\begin{axis}[width=0.47\linewidth, height=5.0cm, axis lines=middle,
  xlabel={$U$}, ylabel={$I$}, xmin=-2.6, xmax=4.6, ymin=0, ymax=1.25,
  xtick={-1.8}, xticklabels={$-U_{\text з}$}, ytick=\empty,
  xlabel style={at={(ticklabel* cs:1)}, anchor=west},
  ylabel style={at={(ticklabel* cs:1)}, anchor=south}, clip=false]
  \addplot[domain=-1.8:4.5, samples=80, thick, accent]
     {1.0*(1-exp(-1.15*(x+1.8)))};
  \addplot[domain=-1.8:4.5, samples=80, thick, keycol, dashed]
     {0.55*(1-exp(-1.15*(x+1.8)))};
  \node[font=\small, accent] at (axis cs:3.2,1.1) {$I_1$};
  \node[font=\small, keycol] at (axis cs:3.2,0.66) {$I_2<I_1$};
\end{axis}
\end{tikzpicture}\hfill
\begin{tikzpicture}
\begin{axis}[width=0.47\linewidth, height=5.0cm, axis lines=middle,
  xlabel={$\nu$}, ylabel={$U_{\text з}$}, xmin=0, xmax=10, ymin=-2.2, ymax=4.2,
  xtick={4}, xticklabels={$\nu_0$}, ytick=\empty,
  xlabel style={at={(ticklabel* cs:1)}, anchor=west},
  ylabel style={at={(ticklabel* cs:1)}, anchor=south}, clip=false]
  \addplot[domain=4:9.5, thick, accent] {0.65*(x-4)};
  \addplot[domain=0:4, thin, dashed, gray] {0.65*(x-4)};
  \node[font=\small, anchor=west] at (axis cs:6.3,1.1) {$\tg\varphi = h/e$};
  \node[font=\small, anchor=west] at (axis cs:0.2,-2.0) {$-A/e$};
\end{axis}
\end{tikzpicture}
\caption{Слева: вольт-амперные характеристики фотоэлемента при одной и той же частоте и
двух интенсивностях света. Ток насыщения пропорционален интенсивности, задерживающее
напряжение $U_{\text з}$ одинаково. Справа: зависимость $U_{\text з}$ от частоты — прямая с
наклоном $h/e$, отсекающая на оси частот красную границу $\nu_0$ (качественно).}
\label{fig:1:photo}
\end{figure}

\paragraph{Почему не работает волновая картина.} В классической теории энергия, поглощаемая
электроном, пропорциональна интенсивности волны и времени облучения; поэтому $T_{\max}$ должна
расти с интенсивностью, красной границы быть не должно, а при слабом свете электрон должен
долго «накапливать» энергию. Оценка: пусть свет с интенсивностью $I=10^{-6}$~Вт/см$^2$
поглощается площадкой атомного размера $\sigma\sim10^{-15}$~см$^2$. Поглощаемая мощность
$I\sigma\sim10^{-21}$~Вт $\approx 6\cdot10^{-3}$~эВ/с, и для набора энергии $\sim2$~эВ
потребовалось бы $\sim 3\cdot 10^{2}$~с — в противоречии с безынерционностью фотоэффекта.
(Оценка предполагает, что электрон поглощает энергию лишь с площадки атомного размера.)

\subsection{Уравнение Эйнштейна}

\begin{principle}{Уравнение Эйнштейна для фотоэффекта}
Фотон поглощается одним электроном целиком. Энергия фотона расходуется на работу выхода
и сообщение электрону кинетической энергии:
\begin{equation}\label{eq:1:einstein}
  \hbar\omega = A + T_{\max},\qquad T_{\max}=\frac{\mel v_{\max}^2}{2}=eU_{\text з}.
\end{equation}
\end{principle}

Уравнение~\eqref{eq:1:einstein} объясняет все законы фотоэффекта.
(1) Число фотоэлектронов в единицу времени пропорционально числу поглощённых фотонов, т.\,е.
интенсивности — отсюда первый закон. (2) $T_{\max}$ определяется энергией одного фотона, а не
их числом. (3) При $\hbar\omega<A$ ни один фотон не может освободить электрон; многократное
поглощение при обычных интенсивностях маловероятно — отсюда красная граница. (4) Акт
поглощения фотона мгновенный, накапливать энергию не требуется.

Знак «max»: электроны, находящиеся глубже уровня Ферми или теряющие энергию при выходе из
металла, вылетают с меньшей энергией; уравнение~\eqref{eq:1:einstein} описывает электроны
с верхнего заполненного уровня, вылетевшие без потерь. Р.~Милликен (1916) проверил линейную
зависимость $U_{\text з}(\nu)$ и по её наклону $h/e$ определил постоянную Планка с точностью около
0{,}5\,\%. Характерные работы выхода: Cs — около 2~эВ ($\lambda_0\approx 0{,}6$~мкм, видимый
свет), Zn — около 4{,}3~эВ, Pt — около 5{,}6~эВ (ультрафиолет).

\begin{remark}
На опыте измеряется разность потенциалов между электродами, к которой добавляется контактная
разность потенциалов катод--анод. Она сдвигает прямую $U_{\text з}(\nu)$, но не меняет её
наклона $h/e$. При очень больших интенсивностях (лазеры) возможен многофотонный фотоэффект:
$N\hbar\omega=A+T_{\max}$, и красная граница в обычном смысле исчезает.
\end{remark}

\paragraph{Роль связанности электрона.} Свободный электрон не может поглотить фотон
(вывод~\ref{der:1:free}): законы сохранения энергии и импульса несовместимы. В фотоэффекте
избыточный импульс передаётся кристаллической решётке (или ядру атома в атомном фотоэффекте),
масса которых велика, поэтому энергию она практически не уносит. Именно поэтому фотоэффект
на атомах наиболее вероятен для электронов внутренних, сильно связанных оболочек.

\subsection{Эффект Комптона}

\paragraph{Опыт.} А.~Комптон (1923) рассеивал рентгеновское излучение (линия K$_\alpha$ молибдена,
$\lambda=0{,}071$~нм) на графите и измерял спектр рассеянного излучения кристаллическим
спектрометром при разных углах рассеяния $\theta$ (рис.~\ref{fig:1:compton}). Наблюдалось:
\begin{itemize}
 \item в рассеянном излучении присутствуют две линии — несмещённая ($\lambda$) и смещённая
       ($\lambda'>\lambda$);
 \item сдвиг $\Delta\lambda=\lambda'-\lambda$ растёт с углом рассеяния и не зависит ни от
       длины волны падающего излучения, ни от вещества рассеивателя;
 \item с ростом атомного номера рассеивателя интенсивность смещённой линии падает, а
       несмещённой растёт.
\end{itemize}
Классическая теория (томсоновское рассеяние) предсказывает, что электрон колеблется в поле волны
с её частотой и излучает на той же частоте, т.\,е. $\lambda'=\lambda$. Смещённая линия
классически необъяснима.

\paragraph{Объяснение.} Энергия рентгеновского фотона ($\hbar\omega\approx17$~кэВ) много больше
энергии связи внешних электронов лёгких атомов (единицы--десятки эВ), поэтому такие электроны
можно считать свободными и покоящимися. Рассеяние — упругое столкновение двух частиц:
фотона с импульсом $\hbar\vect k$ и электрона. Фотон отдаёт часть энергии электрону отдачи,
его частота уменьшается. Из законов сохранения (вывод~\ref{der:1:compton}):
\begin{equation}\label{eq:1:compton}
  \Delta\lambda=\lambda'-\lambda=\lC\,(1-\cos\theta)=2\lC\sin^2\frac{\theta}{2},\qquad
  \lC=\frac{h}{\mel c}.
\end{equation}
Сдвиг меняется от 0 (рассеяние вперёд) до $2\lC\approx4{,}85$~пм (рассеяние назад).
Для видимого света ($\lambda\sim500$~нм) относительный сдвиг $\Delta\lambda/\lambda\sim10^{-5}$,
поэтому эффект заметен лишь в рентгеновском и гамма-диапазонах.

\paragraph{Несмещённая линия.} Фотон может рассеяться на сильно связанном (внутреннем)
электроне так, что импульс отдачи получает атом как целое. В формулу~\eqref{eq:1:compton}
тогда входит масса атома $M\gg\mel$, и сдвиг $h/Mc$ ничтожен. Чем тяжелее атом, тем больше
доля сильно связанных электронов и тем интенсивнее несмещённая линия.

\begin{figure}[htbp]
\centering
\begin{tikzpicture}[scale=1.0]
  \draw[thick, decorate, decoration={snake, amplitude=1.2mm, segment length=3mm}, ->]
     (-3.6,0) -- (-0.25,0);
  \node[above] at (-2.0,0.25) {$\hbar\omega,\ \hbar\vect k$};
  \fill[accent] (0,0) circle (2.4pt);
  \node[below left] at (-0.1,-0.1) {$e^-$ (покоится)};
  \draw[thick, decorate, decoration={snake, amplitude=1.2mm, segment length=3mm}, ->]
     (0.2,0.12) -- (3.0,1.95);
  \node[above left] at (2.3,1.6) {$\hbar\omega',\ \hbar\vect k'$};
  \draw[very thick, keycol, ->] (0.2,-0.1) -- (2.6,-1.6);
  \node[right] at (2.6,-1.6) {$\vect p_e,\ E_e$};
  \draw[dashed, gray] (0,0) -- (3.6,0);
  \draw[->] (1.3,0) arc[start angle=0, end angle=33, radius=1.3];
  \node at (1.6,0.42) {$\theta$};
  \draw[->] (1.1,0) arc[start angle=0, end angle=-32, radius=1.1];
  \node at (1.38,-0.36) {$\varphi$};
  % треугольник импульсов
  \begin{scope}[xshift=6.2cm, yshift=-0.6cm]
    \draw[very thick, ->] (0,0) -- (3.2,0) node[midway, below] {$\hbar\vect k$};
    \draw[very thick, accent, ->] (0,0) -- (2.0,1.4) node[midway, above left] {$\hbar\vect k'$};
    \draw[very thick, keycol, ->] (2.0,1.4) -- (3.2,0) node[midway, above right] {$\vect p_e$};
    \draw (0.7,0) arc[start angle=0, end angle=35, radius=0.7];
    \node at (0.95,0.25) {$\theta$};
  \end{scope}
\end{tikzpicture}
\caption{Геометрия эффекта Комптона (слева) и треугольник импульсов
$\hbar\vect k=\hbar\vect k'+\vect p_e$ (справа). $\theta$ — угол рассеяния фотона,
$\varphi$ — угол вылета электрона отдачи.}
\label{fig:1:compton}
\end{figure}

\section{Математические выводы}

\begin{derivation}[label=der:1:p]{импульс фотона}
Исходное соотношение — релятивистская связь энергии и импульса~\eqref{eq:1:Ep} и формула
для скорости $\vect v=\vect pc^2/E$. Требование $|\vect v|=c$ даёт $pc^2/E=c$, т.\,е.
$E=pc$. Подставив в~\eqref{eq:1:Ep}, получаем $m^2c^4=E^2-p^2c^2=0$. Тогда при $E=\hbar\omega$
\[
  p=\frac{E}{c}=\frac{\hbar\omega}{c}=\hbar k=\frac{h}{\lambda}.
\]
Размерность: $[\hbar k]=\text{эрг}\cdot\text{с}\cdot\text{см}^{-1}=\text{г}\cdot\text{см}/\text{с}$ —
размерность импульса.
\end{derivation}

\begin{derivation}[label=der:1:einst]{уравнение Эйнштейна и задерживающее напряжение}
\emph{Допущения:} фотон поглощается одним электроном; энергия связи наиболее слабо связанного
электрона металла равна работе выхода $A$; потерь при выходе нет.

Закон сохранения энергии для акта поглощения: $\hbar\omega = A + T_{\max}$.
Электрон, вылетевший с энергией $T_{\max}$, в тормозящем поле между катодом и анодом
теряет энергию $eU$. Ток прекращается, когда даже самые быстрые электроны не достигают анода:
\[
  T_{\max}=eU_{\text з}\quad\Longrightarrow\quad
  U_{\text з}=\frac{\hbar}{e}\,\omega-\frac{A}{e}=\frac{h}{e}\,(\nu-\nu_0),\qquad
  \nu_0=\frac{A}{h}.
\]
Наклон прямой $U_{\text з}(\nu)$ равен $h/e$ и не зависит от материала катода, а точка
пересечения с осью частот даёт красную границу. Длина волны красной границы
$\lambda_0=c/\nu_0=hc/A$; в удобных единицах $\lambda_0[\text{Å}]=12\,398/A[\text{эВ}]$.
\end{derivation}

\begin{derivation}[label=der:1:free]{свободный электрон не может поглотить фотон}
Перейдём в систему отсчёта, где электрон до поглощения покоится. До поглощения: энергия
$\hbar\omega+\mel c^2$, импульс $\hbar k$. После поглощения есть только электрон с энергией $E$ и
импульсом $p$. Законы сохранения:
\[
  \hbar\omega+\mel c^2=E,\qquad \hbar k=\frac{\hbar\omega}{c}=p .
\]
Подставим в релятивистское соотношение $E^2=p^2c^2+\mel^2c^4$:
\begin{align*}
  (\hbar\omega+\mel c^2)^2 &= (\hbar\omega)^2+\mel^2c^4,\\
  (\hbar\omega)^2+2\hbar\omega\,\mel c^2+\mel^2c^4 &= (\hbar\omega)^2+\mel^2c^4,\\
  2\hbar\omega\,\mel c^2 &=0 \quad\Longrightarrow\quad \omega=0 .
\end{align*}
Противоречие означает, что поглощение фотона свободным электроном невозможно.
Для фотоэффекта необходима «третья частица» (решётка, ядро), принимающая импульс.
\end{derivation}

\begin{derivation}[label=der:1:compton]{формула Комптона}
\emph{Допущения:} электрон свободен и до столкновения покоится; столкновение упругое;
используются релятивистские формулы (электрон отдачи может быть релятивистским).

\emph{Законы сохранения} энергии и импульса:
\begin{align}
  \hbar\omega+\mel c^2 &= \hbar\omega'+E_e, \label{eq:1:cE}\\
  \hbar\vect k &= \hbar\vect k'+\vect p_e. \label{eq:1:cp}
\end{align}
\emph{Шаг 1.} Из~\eqref{eq:1:cp} $\vect p_e=\hbar(\vect k-\vect k')$; возведём в квадрат,
учитывая, что угол между $\vect k$ и $\vect k'$ равен $\theta$ и $k=\omega/c$:
\[
  p_e^2c^2=\hbar^2\bigl(\omega^2+\omega'^2-2\omega\omega'\cos\theta\bigr).
\]
\emph{Шаг 2.} Из~\eqref{eq:1:cE} $E_e=\hbar(\omega-\omega')+\mel c^2$; возведём в квадрат:
\[
  E_e^2=\hbar^2(\omega-\omega')^2+2\hbar(\omega-\omega')\mel c^2+\mel^2c^4 .
\]
\emph{Шаг 3.} Подставим оба выражения в $E_e^2=p_e^2c^2+\mel^2c^4$:
\begin{align*}
  \hbar^2(\omega^2-2\omega\omega'+\omega'^2)+2\hbar(\omega-\omega')\mel c^2
     &=\hbar^2(\omega^2+\omega'^2-2\omega\omega'\cos\theta),\\
  2\hbar(\omega-\omega')\mel c^2 &= 2\hbar^2\omega\omega'(1-\cos\theta),\\
  \mel c^2\Bigl(\frac{1}{\omega'}-\frac{1}{\omega}\Bigr) &= \hbar(1-\cos\theta)
  \qquad\bigl(\text{деление на }2\hbar\,\omega\omega'\bigr).
\end{align*}
\emph{Шаг 4.} Переходя к длинам волн, $1/\omega=\lambda/(2\pi c)$:
\[
  \frac{\mel c^2}{2\pi c}(\lambda'-\lambda)=\hbar(1-\cos\theta)
  \quad\Longrightarrow\quad
  \boxed{\;\lambda'-\lambda=\frac{2\pi\hbar}{\mel c}(1-\cos\theta)=\lC(1-\cos\theta)\;}
\]
\emph{Проверка размерности:} $[h/(mc)]=\dfrac{\text{г}\cdot\text{см}^2/\text{с}}{\text{г}\cdot\text{см}/\text{с}}=\text{см}$.
\emph{Предельные случаи:} при $\theta=0$ сдвига нет (фотон не передаёт импульс); при
$\theta=\pi$ сдвиг максимален, $2\lC$.
\end{derivation}

\begin{derivation}[label=der:1:recoil]{энергии рассеянного фотона и электрона отдачи, угол отдачи}
Обозначим $\varepsilon=\hbar\omega/\mel c^2$. Из последней строки шага 3 вывода~\ref{der:1:compton}
$1/\omega'=1/\omega+\hbar(1-\cos\theta)/\mel c^2$, откуда
\begin{equation}\label{eq:1:omegap}
  \hbar\omega'=\frac{\hbar\omega}{1+\varepsilon(1-\cos\theta)} .
\end{equation}
Кинетическая энергия электрона отдачи $T_e=\hbar\omega-\hbar\omega'$:
\begin{equation}\label{eq:1:Te}
  T_e=\hbar\omega\,\frac{\varepsilon(1-\cos\theta)}{1+\varepsilon(1-\cos\theta)},\qquad
  T_e^{\max}=T_e(\theta=\pi)=\hbar\omega\,\frac{2\varepsilon}{1+2\varepsilon}.
\end{equation}
Угол вылета электрона $\varphi$ найдём из проекций~\eqref{eq:1:cp} на направление $\vect k$ и
перпендикулярное ему:
\[
  \hbar k=\hbar k'\cos\theta+p_e\cos\varphi,\qquad 0=\hbar k'\sin\theta-p_e\sin\varphi
  \;\Longrightarrow\;
  \tg\varphi=\frac{k'\sin\theta}{k-k'\cos\theta}.
\]
По~\eqref{eq:1:omegap} $k/k'=1+\varepsilon(1-\cos\theta)$, поэтому
\[
  \tg\varphi=\frac{\sin\theta}{1+\varepsilon(1-\cos\theta)-\cos\theta}
  =\frac{\sin\theta}{(1+\varepsilon)(1-\cos\theta)}
  =\frac{\ctg(\theta/2)}{1+\varepsilon},
\]
т.\,е. $\ctg\varphi=(1+\varepsilon)\tg(\theta/2)$.
Электрон отдачи всегда летит вперёд ($0\le\varphi\le\pi/2$).
\end{derivation}

\section{Итоговые формулы}

\keyformula{E=\hbar\omega,\qquad \vect p=\hbar\vect k,\qquad E=pc}
 {$\omega$ — круговая частота, $\vect k$ — волновой вектор ($k=2\pi/\lambda$), $\hbar=h/2\pi$.}
 {всегда для фотона в вакууме ($m=0$).}
 {связывает волновые характеристики излучения с корпускулярными.}

\keyformula{\hbar\omega=A+\frac{\mel v_{\max}^2}{2},\qquad eU_{\text з}=\hbar\omega-A,\qquad \omega_0=\frac{A}{\hbar}}
 {$A$ — работа выхода, $v_{\max}$ — максимальная скорость фотоэлектронов, $U_{\text з}$ — задерживающее напряжение.}
 {однофотонный фотоэффект, нерелятивистские фотоэлектроны.}
 {найти $T_{\max}$, красную границу, постоянную $h/e$ по наклону $U_{\text з}(\nu)$.}

\keyformula{\Delta\lambda=\lC(1-\cos\theta)=2\lC\sin^2\frac{\theta}{2},\qquad \lC=\frac{h}{\mel c}=2{,}43\pmet}
 {$\theta$ — угол рассеяния фотона, $\mel$ — масса электрона.}
 {рассеяние на свободном покоящемся электроне ($\hbar\omega\gg$ энергии связи).}
 {сдвиг длины волны; не зависит от $\lambda$ и вещества.}

\keyformula{\hbar\omega'=\frac{\hbar\omega}{1+\dfrac{\hbar\omega}{\mel c^2}(1-\cos\theta)},\qquad
  T_e^{\max}=\hbar\omega\frac{2\varepsilon}{1+2\varepsilon},\ \ \varepsilon=\frac{\hbar\omega}{\mel c^2}}
 {$\hbar\omega,\hbar\omega'$ — энергии падающего и рассеянного фотонов; $T_e$ — кинетическая энергия электрона отдачи.}
 {те же, что для формулы Комптона.}
 {энергия рассеянного кванта и максимальная энергия электрона отдачи («комптоновский край»).}

\section{Физическая интерпретация}

\textbf{Предельные случаи.} При $\hbar\omega\ll\mel c^2$ ($\varepsilon\ll1$) энергия фотона
почти не меняется, $\hbar\omega'\approx\hbar\omega$, — рассеяние переходит в классическое
томсоновское. При $\varepsilon\gg1$ фотон, рассеянный назад, имеет энергию $\hbar\omega'\to
\mel c^2/2\approx 0{,}26$~МэВ независимо от исходной: почти вся энергия уходит электрону.

\textbf{Смысл комптоновской длины волны.} $\lC$ — длина волны фотона с энергией $\mel c^2$
($\hbar\omega=\mel c^2\iff\lambda=h/\mel c$). Она задаёт масштаб, на котором проявляются
релятивистские квантовые эффекты: попытка локализовать электрон в области размером
$\Delta x\lesssim\hbar/\mel c$ требует импульсов $\Delta p\gtrsim\mel c$ и энергий $\gtrsim\mel c^2$,
достаточных для рождения электрон-позитронных пар. Иерархия атомных масштабов:
$r_e=\alpha\bar\lambda_{\mathrm C}$ (классический радиус), $\bar\lambda_{\mathrm C}$,
$\aB=\bar\lambda_{\mathrm C}/\alpha$, где $\alpha\approx1/137$.

\textbf{Связь с экспериментом.} Независимость $\Delta\lambda$ от вещества указывает, что
рассеяние идёт на электронах — частицах, одинаковых во всех веществах. Совпадение измеренного
$\Delta\lambda(\theta)$ с~\eqref{eq:1:compton} и регистрация электрона отдачи в совпадении с
рассеянным фотоном (опыты Боте и Гейгера, Комптона и Саймона) подтвердили, что энергия и
импульс сохраняются в каждом элементарном акте, а не только в среднем.

\textbf{Двойственность.} Ни волновая, ни корпускулярная модель по отдельности не описывает свет
полностью: в распространении и интерференции проявляются волновые свойства, а в обмене
энергией и импульсом с веществом — корпускулярные.

\section{Что необходимо знать на экзамене}

\begin{examcheck}
\checkhead{Определения}
\begin{itemize}
 \item фотон: $E=\hbar\omega$, $\vect p=\hbar\vect k$, $m=0$, два состояния поляризации;
 \item работа выхода, красная граница, задерживающее напряжение, ток насыщения;
 \item эффект Комптона; комптоновская длина волны $\lC=h/\mel c\approx2{,}43$~пм.
\end{itemize}
\checkhead{Формулы}
\begin{itemize}
 \item $\hbar\omega=A+T_{\max}$, $eU_{\text з}=\hbar\omega-A$, $\lambda_0=hc/A$;
 \item $\Delta\lambda=\lC(1-\cos\theta)$; $\hbar\omega'=\hbar\omega/[1+\varepsilon(1-\cos\theta)]$.
\end{itemize}
\checkhead{Выводы, которые нужно уметь воспроизвести}
\begin{itemize}
 \item $E=pc$ для частицы с $m=0$;
 \item невозможность поглощения фотона свободным электроном;
 \item формула Комптона из законов сохранения (шаги 1--4 вывода~\ref{der:1:compton}).
\end{itemize}
\checkhead{Интерпретации}
\begin{itemize}
 \item почему волновая теория не объясняет фотоэффект (интенсивность, красная граница, безынерционность);
 \item происхождение несмещённой линии; почему эффект не виден для видимого света;
 \item физический смысл $\lC$ как масштаба релятивистской локализации.
\end{itemize}
\end{examcheck}
